% Kleine maten
% Not a real \chapter (it's a hand-styled heading like Register's), so it
% needs its own Inhoud/PDF-bookmark entry, same as imakeidx's intoc=true
% does for Register.
\phantomsection\addcontentsline{toc}{chapter}{Kleine maten}
\vspace*{0.4cm}
{\handfont\fontsize{60}{58}\selectfont\color{terracotta}Kleine\\maten\par}
\vspace{8pt}{\color{terracotta}\rule{2.6cm}{0.5pt}\;\raisebox{-0.2ex}{$\scriptstyle\blacklozenge$}}\par
\vspace{1.4em}

\lettrine{E}{en} kookboek gebruikt weleens een maat zonder vaste
hoeveelheid: een scheutje water, een snufje zout, een klontje boter. Die
termen zijn expres vaag, maar hebben wel een ongeveer vaste orde van
grootte. Hieronder staan ze op een rijtje: van klein naar groot langs een
lijn, met een korte omschrijving en een indicatie in gram. Het zijn
algemene kookschattingen om de termen naast elkaar te kunnen zetten, geen
exacte maten uit dit boek.

% Verticale tijdlijn: één middenlijn met een terracotta ruitje per term, en
% de naam plus een korte omschrijving afwisselend links en rechts ervan. De
% afstand tussen de termen staat niet op schaal (dat zou de tekstblokjes
% van de kleinste termen te dicht op elkaar duwen); de volgorde van klein
% naar groot en het grammage in elk tekstblokje geven de verhouding weer.
\begin{center}
\begin{tikzpicture}[y=1cm]
\draw[color=terracotta, line width=0.6pt] (0,0.3) -- (0,-13.8);

\node at (0,0) {\textcolor{terracotta}{\scriptsize$\blacklozenge$}};
\draw[color=inkfaint, line width=0.3pt] (0,0) -- (-0.5,0);
\node[align=right, text width=4.2cm, anchor=east] at (-0.6,0)
  {\textbf{\textcolor{terracotta}{Mespunt}} \textcolor{inkfaint}{(± 0,3--0,5 g)}\\[1pt]
   {\footnotesize\textcolor{inkfaint}{Het beetje dat op de punt van een mes blijft liggen.}}};

\node at (0,-1.5) {\textcolor{terracotta}{\scriptsize$\blacklozenge$}};
\draw[color=inkfaint, line width=0.3pt] (0,-1.5) -- (0.5,-1.5);
\node[align=left, text width=4.2cm, anchor=west] at (0.6,-1.5)
  {\textbf{\textcolor{terracotta}{Snufje}} \textcolor{inkfaint}{(± 0,5--1 g)}\\[1pt]
   {\footnotesize\textcolor{inkfaint}{Wat er tussen duim en twee vingers past.}}};

\node at (0,-3) {\textcolor{terracotta}{\scriptsize$\blacklozenge$}};
\draw[color=inkfaint, line width=0.3pt] (0,-3) -- (-0.5,-3);
\node[align=right, text width=4.2cm, anchor=east] at (-0.6,-3)
  {\textbf{\textcolor{terracotta}{Takje}} \textcolor{inkfaint}{(± 1--2 g)}\\[1pt]
   {\footnotesize\textcolor{inkfaint}{Eén twijgje van een houtig kruid.}}};

\node at (0,-4.5) {\textcolor{terracotta}{\scriptsize$\blacklozenge$}};
\draw[color=inkfaint, line width=0.3pt] (0,-4.5) -- (0.5,-4.5);
\node[align=left, text width=4.2cm, anchor=west] at (0.6,-4.5)
  {\textbf{\textcolor{terracotta}{Theelepel}} \textcolor{inkfaint}{(± 2--3 g)}\\[1pt]
   {\footnotesize\textcolor{inkfaint}{Van een gemalen specerij.}}};

\node at (0,-6) {\textcolor{terracotta}{\scriptsize$\blacklozenge$}};
\draw[color=inkfaint, line width=0.3pt] (0,-6) -- (-0.5,-6);
\node[align=right, text width=4.2cm, anchor=east] at (-0.6,-6)
  {\textbf{\textcolor{terracotta}{Eetlepel}} \textcolor{inkfaint}{(± 6--8 g)}\\[1pt]
   {\footnotesize\textcolor{inkfaint}{Van een gemalen specerij, afgestreken.}}};

\node at (0,-7.5) {\textcolor{terracotta}{\scriptsize$\blacklozenge$}};
\draw[color=inkfaint, line width=0.3pt] (0,-7.5) -- (0.5,-7.5);
\node[align=left, text width=4.2cm, anchor=west] at (0.6,-7.5)
  {\textbf{\textcolor{terracotta}{Flinke lepel}} \textcolor{inkfaint}{(± 10--12 g)}\\[1pt]
   {\footnotesize\textcolor{inkfaint}{Een eetlepel, maar dan gehoopt.}}};

\node at (0,-9) {\textcolor{terracotta}{\scriptsize$\blacklozenge$}};
\draw[color=inkfaint, line width=0.3pt] (0,-9) -- (-0.5,-9);
\node[align=right, text width=4.2cm, anchor=east] at (-0.6,-9)
  {\textbf{\textcolor{terracotta}{Klontje boter}} \textcolor{inkfaint}{(± 10--15 g)}\\[1pt]
   {\footnotesize\textcolor{inkfaint}{Ongeveer zo groot als een dobbelsteen.}}};

\node at (0,-10.5) {\textcolor{terracotta}{\scriptsize$\blacklozenge$}};
\draw[color=inkfaint, line width=0.3pt] (0,-10.5) -- (0.5,-10.5);
\node[align=left, text width=4.2cm, anchor=west] at (0.6,-10.5)
  {\textbf{\textcolor{terracotta}{Bosje kruiden}} \textcolor{inkfaint}{(± 15--30 g)}\\[1pt]
   {\footnotesize\textcolor{inkfaint}{Een standaardbosje, zoals bij de groenteboer.}}};

\node at (0,-12) {\textcolor{terracotta}{\scriptsize$\blacklozenge$}};
\draw[color=inkfaint, line width=0.3pt] (0,-12) -- (-0.5,-12);
\node[align=right, text width=4.2cm, anchor=east] at (-0.6,-12)
  {\textbf{\textcolor{terracotta}{Handje}} \textcolor{inkfaint}{(± 30--50 g)}\\[1pt]
   {\footnotesize\textcolor{inkfaint}{Wat je losjes in één hand houdt.}}};

\node at (0,-13.5) {\textcolor{terracotta}{\scriptsize$\blacklozenge$}};
\draw[color=inkfaint, line width=0.3pt] (0,-13.5) -- (0.5,-13.5);
\node[align=left, text width=4.2cm, anchor=west] at (0.6,-13.5)
  {\textbf{\textcolor{terracotta}{Scheutje}} \textcolor{inkfaint}{(± 40--60 ml)}\\[1pt]
   {\footnotesize\textcolor{inkfaint}{Bijvoorbeeld water om af te blussen.}}};
\end{tikzpicture}
\end{center}

Sterke gemalen specerijen zoals kerrie, komijn, kaneel, kurkuma,
gemberpoeder en chilipoeder rekenen we daarom in theelepels, niet in
eetlepels: een eetlepel is al snel twee tot drie keer zoveel en overheerst
al gauw de rest van het gerecht.

\vspace{0.6em}
{\footnotesize\itshape\color{inkfaint}Bij twijfel geldt: liever een beetje
minder dan te veel. Je kunt op het eind altijd bijproeven en bijkruiden.\par}
